% Faithful formal-publication transcription of TheoremB' and complete Section5, pp16:6,16:21--28.
\section*{Author main counting statement}
\begin{authorclaim}{TheoremB'}
For $q$ a prime power with $q=n^{\Omega(1)}$, there is a randomized counting-to-decision reduction for Alternating Matrix Space Isometry which, given $n\times n$ alternating matrix spaces $\mathcal A,\mathcal B$ over $\mathbb F_q$ of dimension $m$, computes the number of isometries from $\mathcal A$ to $\mathcal B$ in time $q^{O(n)}$. The reduction queries the decision oracle with inputs of dimension at most $O(n^2)$.
\end{authorclaim}
\section*{5 Counting-to-decision reduction by restricting to diagonal groups}
In this section, we devise a gadget to achieve the restriction to the group of diagonal matrices, and use it to do the counting to decision reduction for Alternating Matrix Space Isometry.
\subsection*{5.1 Preliminaries}
Some preparations are in order.
\begin{authorclaim}{Observation12}
Let $n\ge23$. Then any permutation $\sigma\in S_n$ either fixes a set of 6 points $P\subseteq[n]$, or moves a set of 6 points $P\subseteq[n]$ to another set of 6 points $Q\subseteq[n]$ such that these two sets are disjoint.
\end{authorclaim}
\begin{proof}
Suppose $\sigma$ fixes at most 5 points. Then there are at least 18 points that are not fixed by $\sigma$. Suppose $\sigma$ has $t$ non-trivial cycles of length $l_1,\ldots,l_t$, such that $\sum_i l_i\ge18$. For a cycle $(p_1,\ldots,p_s)$, we can choose $p_1,p_3,\ldots,p_{2\lfloor s/2\rfloor-1}$ and put them in $P$, and $p_2,p_4,\ldots,p_{2\lfloor s/2\rfloor}$ in $Q$. Do this for every cycle, we obtain the desired $P$ and $Q$. The worst case is when every cycle is of length 3. Since there are at least 18 points not fixed by $\sigma$, $P$ is of size $\ge6$.
\end{proof}
We shall make repeated uses of the following facts.
\begin{authorclaim}{Fact13}
\begin{enumerate}
\item Given $a_i\in\mathbb R$, $0\le a_i\le1$, $i\in[m]$, $\prod_{i\in[m]}(1-a_i)\ge1-\sum_{i\in[m]}a_i$.
\item Let $m,N\in\mathbb N$ and $1\le m\le N$. A random matrix $A\in\mathrm M(N\times m,q)$ is of rank $m$ with probability $\ge1-2/q^{N-m+1}$.
\item For $n\in\mathbb N$, $0\le d\le n$, the number of dimension-$d$ subspaces of $\mathbb F_q^n$ is equal to the Gaussian binomial coefficient
\[\binom nd_q:=\frac{(q^n-1)\cdot(q^n-q)\cdots(q^n-q^{d-1})}{(q^d-1)\cdot(q^d-q)\cdots(q^d-q^{d-1})}.\]
\item The Gaussian binomial coefficient satisfies $q^{(n-d)d}\le\binom nd_q\le q^{(n-d)d+d}$.
\item For $d\in\mathbb N$, the number of complement subspaces of a fixed dimension-$d$ subspace of $\mathbb F_q^n$ is $q^{d(n-d)}$.
\end{enumerate}
\end{authorclaim}
\begin{proof}
For(2), $\Pr[\operatorname{rk}(A)=m]=(1-1/q^N)\cdot(1-q/q^N)\cdots(1-q^{m-1}/q^N)$. By(1), we have
\[\Pr[\operatorname{rk}(A)=m]\ge1-\sum_{i=N-m+1}^N\frac1{q^i}=1-\frac1{q^{N-m+1}}-\sum_{i=N-m+2}^N\frac1{q^i}\ge1-\frac2{q^{N-m+1}}.\]
\end{proof}
\subsection*{5.2 Describing the gadget}
Let $\mathcal A\le\Lambda(n,q)$ be an alternating matrix space, and let $\mathbf A=(A_1,\ldots,A_m)\in\Lambda(n,q)^m$ be an ordered linear basis of $\mathcal A$. Let $\mathtt A\in\mathrm T(n\times n\times m,\mathbb F_q)$ be the 3-way array constructed from $\mathbf A$, i.e. the $i$th frontal slice of $\mathtt A$ is $A_i$.
We shall assume $n=\Omega(1)$, and $q=n^{\Omega(1)}$ throughout the remainder of this section.
\paragraph{The form of the gadget.} To describe the gadget, it is easier to view $\mathtt A$ from the lateral viewpoint. That is, for $i\in[n]$, let $C_i=[A_1e_i,\ldots,A_me_i]\in\mathrm M(n\times m,q)$. Let $\mathbf C=(C_1,\ldots,C_n)\in\mathrm M(n\times m,q)^n$. Then construct $\mathbf C'=(C'_1,\ldots,C'_n)$, $C'_i=\begin{bmatrix}C_i&0\\0&G_i\end{bmatrix}$, where $G_i$ is of size $6n\times4n^2$. For $i\in[n]$, $G_i=[0\ \cdots\ 0\ H_i\ 0\ \cdots\ 0]$, where $H_i$ is of size $6n\times4n$ in the $i$th block, and 0 denotes an all-zero matrix of size $6n\times4n$. The $H_i$ will be described below.
Note that from the frontal viewpoint of looking at $\mathtt A$, $G_i$'s are inserted, vertically, below and behind $\mathtt A$. So to preserve the alternating structure, $-G_i$'s also need to be inserted, horizontally, on the right and behind $\mathtt A$. We therefore get $\widetilde{\mathtt A}$, which is of size $7n\times7n\times(m+4n^2)$.
\paragraph{Conditions imposed on the $H_i$'s.} Of course, the key to the construction above lies in the properties of the $H_i$'s. Let $V_i\le\mathbb F_q^{6n}$ be the subspace spanned by the columns of $H_i$. We shall impose the following conditions on $H_i$.
\begin{enumerate}
\item For any $i\in[n]$, $\operatorname{rk}(H_i)=\dim(V_i)=4n$.
\item For any $i,j\in[n]$, $i\ne j$, $\operatorname{rk}([H_i\ H_j])=\dim(V_i\cup V_j)=6n$.
\item For any $(i_1,i_2,i_3,i_4,i_5,i_6)\in[n]^6$ and $(j_1,j_2,j_3,j_4,j_5,j_6)\in[n]^6$, such that $|\{i_1,\ldots,i_6\}\cup\{j_1,\ldots,j_6\}|=12$, i.e. $i_k$ and $j_\ell$ all different, the coset $\mathcal C=\{T\in\mathrm{GL}(6n,q):\forall k\in[6],T(V_{i_k})=V_{j_k}\}$ is empty. Note that for any $i\in[n]$, $T(V_i)$ is spanned by the columns of $TH_i$.
\item For any $(i_1,i_2,i_3,i_4,i_5,i_6)\in[n]^6$, $i_k$ all different, the group $\mathcal S=\{T\in\mathrm{GL}(6n,q):\forall k\in[6],T(V_{i_k})=V_{i_k}\}$ consists of only of scalar matrices.
\end{enumerate}
\begin{authorclaim}{Remark14}
Given $H_1,\ldots,H_n\in\mathrm M(6n\times4n,q)$, whether they satisfy the four conditions can be verified in polynomial time.
\end{authorclaim}
Conditions(1) and(2) are easily verified in deterministic polynomial time.
For condition(3), it can be formulated as a linear algebraic problem as follows. Let $X$ be a $6n\times6n$ variable matrix. Let $Y_k$, $k\in[6]$, be $4n\times4n$ variable matrices. Set up the equations $XH_{i_k}=H_{j_k}Y_k$, and solve the linear equations to get a subspace of $\mathbb F_q^{(6n)^2+6\cdot(4n)^2}$. The question is then whether this subspace contains $(T,R_1,\ldots,R_6)$ where $T\in\mathrm{GL}(6n,q)$ and $R_i\in\mathrm{GL}(4n,q)$. This is an instance of the symbolic determinant identity testing (SDIT) problem, so it admits a randomized efficient algorithm when $q=n^{\Omega(1)}$.
In fact, this instance of SDIT problem can be solved in deterministic polynomial time. For this let us also check out condition(4). Here, let $X$ and $Y_i$ be from above, and set up the equations $XH_{i_k}=H_{i_k}Y_k$. Solve the linear equations to get a subspace of $\mathbb F_q^{(6n)^2+6\cdot(4n)^2}$. This subspace turns out to be an algebra under the natural multiplications. Indeed, if $AH_{i_k}=H_{i_k}B_k$ and $A'H_{i_k}=H_{i_k}B'_k$, then $AA'H_{i_k}=H_{i_k}B_kB'_k$. To compute the unit group in a matrix algebra can be solved by a polynomial-time Las Vegas algorithm by [16]. Given the unit group, whether it consists of only scalar matrices can be verified easily in deterministic polynomial time.
Then the linear space in condition(3) is a module over the algebra defined in the last paragraph. Because of this structure, the SDIT problem for such instances can be solved in deterministic polynomial time [14,19,37].
\subsection*{5.3 Construction and properties of the gadget}
The following three propositions reveal the construction and functions of the gadget described above. First about the construction. Instead of constructing the above $H_i$'s explicitly in a deterministic way, we shall show that random choices suffice.
\begin{authorclaim}{Proposition15}
Let $H_i\in\mathrm M(6n\times4n,q)$, $i\in[n]$, be random matrices. Then $H_i$'s satisfy the four conditions in Section5.2 with probability $\ge1-n^{O(1)}/q^{\Omega(1)}$.
\end{authorclaim}
Second about the functionality. The following proposition formally explains this.
\begin{authorclaim}{Proposition16}
Suppose $\mathtt A$ and $\mathtt B$ are two 3-tensors constructed from ordered bases of $m$-dimensional alternating matrix spaces $\mathcal A,\mathcal B\le\Lambda(n,q)$. Let $\widetilde{\mathtt A}$ and $\widetilde{\mathtt B}$ be constructed as above, and let $\widetilde{\mathcal A}$ and $\widetilde{\mathcal B}$ be the alternating matrix spaces spanned by the frontal slices of $\widetilde{\mathtt A}$ and $\widetilde{\mathtt B}$, respectively. Then $\mathcal A$ and $\mathcal B$ are isometric via a diagonal matrix if and only if $\widetilde{\mathcal A}$ and $\widetilde{\mathcal B}$ are isometric.
\end{authorclaim}
Finally we shall use this gadget to achieve a counting-to-decision reduction for Alternating Matrix Space Isometry. Formally, we have the following.
\begin{authorclaim}{Proposition17}
Suppose we are given $\mathcal A,\mathcal B\le\Lambda(n,q)$ and a decision oracle for Alternating Matrix Space Isometry. Then there exists a Las Vegas randomized algorithm that computes the number of isometries from $\mathcal A$ to $\mathcal B$ in time $q^{O(n)}$.
\end{authorclaim}
The next three subsections are devoted to the proofs of Propositions15 (Section5.3.3),16 (Section5.3.1), and17 (Section5.3.2). Note that, because the proof of Proposition15 is more complicated compared to the other two, we postpone it to the last.
\begin{authorclaim}{Remark18}
In fact, we expect that this construction works even for small finite fields. The bottleneck lies in Proposition15. If the probability $n^{O(1)}/q^{\Omega(1)}$ could be improved to $n^{O(1)}/q^{\Omega(n)}$, then we would be done. We believe it possible to utilize the structure of invariant subspaces under matrix actions over $\mathbb F_q$ to achieve this. However, we expect that the calculations will be tedious and heavy, so we hope to leave this to a future work.
\end{authorclaim}
\subsection*{5.3.1 Restricting to the diagonal group}
Briefly speaking, conditions 1 and 2 ensure that we first restrict to monomial matrices. Conditions 3 and 4 prevent non-trivial permutations due to the following. As we assume $n=\Omega(1)$, by Observation12, $\sigma\in S_n$ either fixes 6 elements in $[n]$, or moves a set of 6 elements to another, disjoint, set of 6 elements. Condition 3 ensures that the second case could not happen. Condition 4 ensures that in the first case, the only possible invertible matrices that ``preserves'' the matrices $G_i$ for $i\in P$ when multiplying from the left are scalar matrices.
We now prove Proposition16.
\begin{proof}[Proof of Proposition16]
Recall that we construct such $\widetilde{\mathtt A}$ and $\widetilde{\mathtt B}$ from $\mathtt A$ and $\mathtt B$, respectively, using the method in Section5.2. Let $\widetilde{\mathcal A}$ and $\widetilde{\mathcal B}$ be alternating matrix spaces in $\Lambda(7n,q)$, spanned by the frontal slices of $\widetilde{\mathtt A}$ and $\widetilde{\mathtt B}$, respectively.
We want to show that $\widetilde{\mathcal A}$ and $\widetilde{\mathcal B}$ are isometric if and only if $\mathcal A$ and $\mathcal B$ are isometric via diagonal matrices. The if direction is straightforward. Suppose there exist $P=\operatorname{diag}(\alpha_1,\ldots,\alpha_n)\in\operatorname{diag}(n,q)$ and $Q\in\mathrm{GL}(m,q)$ such that $P^t\mathbf A P=\mathbf B^Q$. Let $\widetilde P=\begin{bmatrix}P&0\\0&I_{6n}\end{bmatrix}\in\mathrm{GL}(7n,q)$. Let $\widetilde Q=\begin{bmatrix}Q&0\\0&Q'\end{bmatrix}\in\mathrm{GL}(m+4n^2)$, where $Q'=\operatorname{diag}(\alpha_1I_{4n},\ldots,\alpha_nI_{4n})$. Then it is easy to verify that $\widetilde P^t\widetilde{\mathtt A}\widetilde P=\widetilde{\mathtt B}^{\widetilde Q}$.
Now we turn to the only if direction. If $\widetilde{\mathcal A}$ and $\widetilde{\mathcal B}$ are isometric, then there exists $\widetilde P\in\mathrm{GL}(7n,q)$ and $\widetilde Q\in\mathrm{GL}(m+4n^2,q)$, such that $\widetilde P^t\widetilde{\mathtt A}\widetilde P=\widetilde{\mathtt B}^{\widetilde Q}$. Let $\widetilde P=\begin{bmatrix}P_{1,1}&P_{1,2}\\P_{2,1}&P_{2,2}\end{bmatrix}$, where $P_{1,1}$ is of size $n\times n$. It can be checked easily, from the lateral viewpoint, that $P_{1,2}=0$. As if not, then some $H_i$ would appear in one of the last $6n$ lateral slices in $\widetilde{\mathtt A}^{\widetilde P}$. This would set this slice to be of rank $\ge4n$ by condition(1), which contradicts that the corresponding lateral slice of $\widetilde{\mathtt B}^{\widetilde Q}$ is of rank $\le n$. It follows that $P_{1,1}\in\mathrm{GL}(n,q)$ and $P_{2,2}\in\mathrm{GL}(6n,q)$.
We first claim that $P_{1,1}$ has to be a monomial matrix. If not, then one of the first $n$ lateral slice of $\widetilde{\mathtt A}^{\widetilde P}$ has two distinct $H_i$ and $H_j$. By condition(2), this slice is of rank $\ge6n$, which contradicts that the corresponding lateral slice of $\widetilde{\mathtt B}^{\widetilde Q}$ is of rank $\le5n$.
We further claim that $P_{1,1}$ has to be a diagonal matrix. If not, then suppose the non-trivial permutation underlying $P_{1,1}$ is $\sigma\in S_n$. Since we assumed $n=\Omega(1)$, by Observation12, one of the following two cases has to happen.
\begin{itemize}
\item $\exists\{i_1,\ldots,i_6\}\subseteq[n]$, $\{j_1,\ldots,j_6\}\subseteq[n]$, $|\{i_1,\ldots,i_6\}\cup\{j_1,\ldots,j_6\}|=12$, such that $\sigma(i_k)=j_k$ for $k\in[6]$. We then claim the following.
\begin{authorclaim}{Claim19}
For $\widetilde P^t\widetilde{\mathtt A}\widetilde P=\widetilde{\mathtt B}^{\widetilde Q}$ to hold, a necessary condition is that $\forall k\in[6]$, $P_{2,2}H_{j_k}$ and $H_{i_k}$ have the same linear span.
\end{authorclaim}
\begin{proof}
To see this, note that the $i_k$th lateral slice of $\widetilde P^t\widetilde{\mathtt A}\widetilde P$ is the $j_k$th lateral slice of $\widetilde P^t\widetilde{\mathtt A}$ (up to a scalar multiple). It is equal to the $i_k$th lateral slice of $\widetilde{\mathtt B}^{\widetilde Q}$. Then $\widetilde P^t$ acts on the left on the $j_k$th lateral slice of $\widetilde{\mathtt A}$. Noting that $P^t=\begin{bmatrix}P_{1,1}^t&P_{2,1}^t\\0&P_{2,2}^t\end{bmatrix}$ and the $j_k$th lateral slice of $\widetilde{\mathtt A}$ is $C'_{j_k}=\begin{bmatrix}C_{j_k}&0\\0&G_{j_k}\end{bmatrix}$, we see that $P^tC'_{j_k}=\begin{bmatrix}*&*\\0&P_{2,2}^tG_{j_k}\end{bmatrix}$. (Here, $C_i$ and $G_i$ are defined in Section5.2.) On the other hand, we see that the $i_k$th lateral slice of $\widetilde{\mathtt B}^{\widetilde Q}$ is the $i_k$th lateral slice multiplied from the right by $\widetilde Q$. Our claim follows then by comparing the last $6n$ rows.
\end{proof}
But the condition(3) excludes the existence of such $P_{2,2}$, so this cannot happen.
\item $\exists\{i_1,\ldots,i_6\}\subseteq[n]$, $i_k$ all different, such that $\sigma(i_k)=i_k$. In this case, for $\widetilde P^t\widetilde{\mathtt A}\widetilde P=\widetilde{\mathtt B}^{\widetilde Q}$ to hold, by the same argument as in the proof of Claim19, a necessary condition is that $P_{2,2}H_{i_k}$ and $H_{i_k}$ have the same linear span. Then the condition(4) ensures that $P_{2,2}=\lambda I_{6n}$ for some $\lambda\ne0\in\mathbb F$ in this setting. Then because $\sigma$ is non-trivial, $\sigma$ moves some $i\in[n]$ to $j\in[n]$, $i\ne j$. By comparing the $j$th lateral slice of $\widetilde P^t\widetilde{\mathtt A}$ and the $i$th lateral slice of $\widetilde{\mathtt B}^{\widetilde Q}$, $P_{2,2}H_i=\lambda H_i$ and $H_j$ have the same linear span, which is not possible because the condition(2) ensures that $H_i$ and $H_j$ span different subspaces.
\end{itemize}
We then have shown that $P_{1,1}$ must be a diagonal matrix. By comparing the top-left-front sub-tensors of size $n\times n\times m$ of $\widetilde P^t\widetilde{\mathtt A}\widetilde P$ and $\widetilde{\mathtt B}^{\widetilde Q}$, we arrive at the desired conclusion that $\mathcal A$ and $\mathcal B$ are isometric via the diagonal matrix $P_{1,1}$.
\end{proof}

\subsection*{5.3.2 Using the gadget for counting-to-decision reduction}
The strategy follows closely the counting to decision reduction for graph isomorphism. We then apply the above strategy to get a counting to decision reduction for alternating matrix space isometry to prove Proposition17.
\begin{proof}[Proof of Proposition17]
Our goal is to compute the number of isomorphisms from $\mathcal A$ to $\mathcal B$, where $\mathcal A,\mathcal B\le\Lambda(n,q)$ are of dimension $m$. First, we use the decision oracle first to decide whether $\mathcal A$ and $\mathcal B$ are isometric. If not, the number of isometries is 0. If so, we need to caculate the order of the autometry group of $\mathcal A$, $\operatorname{Aut}(\mathcal A)$. To do that, we first randomly sample $n$ $6n\times4n$ matrices $H_1,\ldots,H_n$ over $\mathbb F_q$, and verify whether they satisfy the four conditions in Section5.2 using Remark14. Note that this is where the algorithm needs to be a Las Vegas algorithm.
Let $A=\operatorname{Aut}(\mathcal A)$. Recall that $e_i$ denotes the $i$th standard basis vector in $\mathbb F_q^n$. For $i\in[n]$, let $A_i=\{T\in A:\forall1\le j\le i,T(e_i)=\lambda_ie_i,\lambda_i\ne0\in\mathbb F_q\}$. Note that $A_n=A\cap\operatorname{diag}(n,q)$. We can calculate the order of $A_n$ in time $q^{O(n)}$ by brute-force, i.e., enumerating all invertible diagonal matrices. Set $A_0=A$. We then have the tower of subgroups $A_0\ge A_1\ge\cdots\ge A_n$.
To compute the order of $A_0$, it is enough to compute $[A_i:A_{i+1}]$. Note that for $T,T'\in A_i$, $TA_{i+1}=T'A_{i+1}$ as left cosets in $A_i$ if and only if $T(e_{i+1})=\lambda T'(e_{i+1})$ for some $\lambda\ne0\in\mathbb F_q$. So $[A_i:A_{i+1}]$ is equal to the size of the orbit of $e_{i+1}$ under $A_i$ in the projective space. Let $v\in\mathbb F_q^n$. To test whether $v$ is in the orbit of $e_{i+1}$ under $A_i$ in the projective space, we tranform $\mathcal A$ by $P^t\cdot P$, where $P\in\mathrm{GL}(n,q)$ sends $e_{i+1}$ to $v$ and $e_j$ to $e_j$ for $j\ne i+1$, to get $\mathcal A'$. We then add the diagonal restriction gadget to the first $i+1$ lateral slices and the first $i+1$ horizontal slices of $\mathcal A$ and $\mathcal A'$, to obtain $\widetilde{\mathcal A}$ and $\widetilde{\mathcal A}'$ respectively. Then feed $\mathcal A$ and $\mathcal A'$ to the decision oracle. By the functionality of the diagonal restriction gadget, $v$ is in the orbit of $e_{i+1}$ in the projective space if and only if $\widetilde{\mathcal A}$ and $\widetilde{\mathcal A}'$ are isometric. Enumerating $v\in\mathbb F_q^n$ up to scalar multiples gives us the size of the orbit of $e_{i+1}$ under $A_i$ in the projective space. This finishes the description of the algorithm.
A small caveat in the above is that our gadget requires $n=\Omega(1)$, so we cannot start from $A_0$ at the beginning. This issue can be revolved by noting that the order of $A_c$, for any constant $c$, can be computed in time $q^{O(n)}$, by enumerating all possible images of $e_1,\ldots,e_c$ in time $q^{O(n)}$, adding the diagonal restriction gadget, and utilizing the decision oracle.
\end{proof}
\subsection*{5.3.3 Random $H_i$'s satisfy the requirements when $q=n^{\Omega(1)}$}
In the following we will encounter random matrices over $\mathbb F_q$ as well as random subspaces in $\mathbb F_q^n$. There is a subtle point which we want to clarify now. Let $m\le n$. Note that there are $\binom nm_q$ subspaces of $\mathbb F_q^n$, and there are $N_1=(q^n-1)\cdots(q^n-q^{m-1})$ rank-$m$ matrices of size $n\times m$. It can be seen easily that each $m$-dimensional subspace $V$ of $\mathbb F_q^n$ has $N_2=(q^m-1)\cdots(q^m-q^{m-1})$ many representations as rank-$m$ matrices of size $n\times m$, i.e. the columns of the matrix span $V$. It follows that we can work with random rank-$m$ matrices of size $n\times m$ as if we are working with random $m$-dimensional subspaces of $\mathbb F_q^n$. Such correspondences will be used implicitly for other structures, including direct sum decompositions.
Now let us get back to our question. We shall show that a random choice of $H_i$, $i\in[n]$, would satisfy the four conditions we imposed on $H_i$'s. We will prove that for conditions $k=1,2,3$,
\[\Pr[\text{random }H_i\text{ not satisfy condition }k]\le\frac{n^{O(1)}}{q^{\Omega(n)}}.\]
Once these hold, by a union bound, we have
\[\Pr[\exists i\in[3],\text{ random }H_i\text{ not satisfy condition }i]\le\frac{n^{O(1)}}{q^{\Omega(n)}}.\]
For condition(4), we will prove that
\[\Pr[\text{random }H_i\text{ not satisfy condition 4}\mid H_i\text{ satisfy conditions 1,2,3}]\le\frac{n^{O(1)}}{q^{\Omega(1)}}.\]
This then would allow us to conclude that when $q=n^{\Omega(1)}$, random $H_i$'s satisfy all the four conditions. We examine the first three conditions one by one.
\begin{enumerate}
\item For condition(1), by Fact13(2), we have $\Pr[\exists i\in[n],\operatorname{rk}(H_i)<4n]\le n\cdot\Pr[\operatorname{rk}(H_i)<4n]\le2n/q^{2n+1}$.
\item For condition(2), noting that the block matrix $(H_i\ H_j)$ is a random $6n\times8n$ matrix over $\mathbb F_q$, by Fact13(2), we have
\[\Pr[\exists i\ne j\in[n],\operatorname{rk}((H_i\ H_j))<6n]\le\binom n2\cdot\frac2{q^{8n-6n+1}}\le\frac{n^2}{q^{2n+1}}.\]
\item For condition(3), let $I=(H_{i_1}\ \cdots\ H_{i_6})$, and $J=(H_{j_1}\ \cdots\ H_{j_6})$. We see that $\mathcal C$ is non-empty if and only if there exists $L\in\mathrm{GL}(6n,q)$ and $R_k\in\mathrm{GL}(4n,q)$, $k\in[6]$, such that $LH_{i_k}R_k=H_{j_k}$. Note that the orbit of $I$ under this group action is of size at most $q^{(6n)^2+6\cdot(4n)^2}=q^{132n^2}$. Since $i_k$ and $j_\ell$ are all different, the probability of $J$ belonging to this orbit is $\le q^{132n^2}/q^{144n^2}=1/q^{12n^2}$. We then have
\[\Pr[\exists i_k,j_k\in[n],k\in[6],i_k,j_k\text{ all different},\mathcal C=\emptyset]\le\binom n{12}\frac2{q^{12n^2}}\le\frac{n^{12}}{q^{12n^2}}.\]
\end{enumerate}
We now focus on condition(4). For condition(4), we first assume that the conditions(1) and(2) as above hold. Then $V_i$'s are random $4n$-dimensional subspaces of $\mathbb F_q^{6n}$. Note that
\[\Pr[\exists i_k\in[n],k\in[6],i_k\text{ all different},\mathcal S\text{ non-scalar}]\le n^6\cdot\Pr[\mathcal S\text{ non-scalar for }V_1,\ldots,V_6].\]
So we turn to study $\Pr[\mathcal S\text{ non-scalar for }V_1,\ldots,V_6]$, and will show that it is $\le1/q^{\Omega(1)}$.
Let $U_1=V_1\cap V_2$, $U_2=V_2\cap V_3$, and $U_3=V_1\cap V_3$. Let $W_1=V_4\cap V_5$, $W_2=V_5\cap V_6$, and $W_3=V_4\cap V_6$. Since conditions(1) and(2) hold, we have $\dim(U_i)=\dim(W_i)=2n$. We claim that with probability $\ge1-2/q$, $\mathbb F_q^{6n}=U_1\oplus U_2\oplus U_3$, i.e., $U_1\cup U_2\cup U_3$ span $\mathbb F_q^{6n}$. This can be seen as follows. Since we assumed conditions(1) and(2), this happens if and only if $V_1\cap V_2$ and $V_3$ together span $\mathbb F_q^{6n}$. Therefore we calculate, using Fact13(1),(3), and(5), that
\begin{align*}
&\Pr[V_3\text{ is a complement subspace of }V_1\cap V_2]\\
&=q^{2n\cdot4n}/\binom{6n}{4n}_q
=\frac{(q^{6n}-q^{2n})(q^{6n}-q^{2n+1})\cdots(q^{6n}-q^{6n-1})}{(q^{6n}-1)(q^{6n}-q)\cdots(q^{6n}-q^{4n-1})}\\
&\ge\frac{(q^{6n}-q^{2n})(q^{6n}-q^{2n+1})\cdots(q^{6n}-q^{6n-1})}{q^{6n}\cdot q^{6n}\cdots q^{6n}}\\
&=(1-1/q^{4n})(1-1/q^{4n-1})\cdots(1-1/q)\ge1-\sum_{i=1}^{4n}1/q^i\ge1-2/q.
\end{align*}
It follows that with probability $\ge1-4/q$, we can assume in addition that $W_i$ form a direct sum decomposition of $\mathbb F_q^{6n}$.
Therefore, we turn to bound the probability that there exists a non-scalar invertible matrix stabilizing these two direct sum decompositions of $\mathbb F_q^{6n}$. Since $i_k$ are all different, the two direct sum decompositions $U_1\oplus U_2\oplus U_3$ and $W_1\oplus W_2\oplus W_3$ are independent.
So we can assume that $U_i$ is spanned by those standard basis vectors $\vec e_{2n(i-1)+1},\ldots,\vec e_{2ni}$, $i=1,2,3$. The group that stabilizes this direct sum decomposition $U_1\oplus U_2\oplus U_3$ consists of $\begin{bmatrix}D_1&0&0\\0&D_2&0\\0&0&D_3\end{bmatrix}\in\mathrm{GL}(6n,\mathbb F_q)$ where $D_i$ is of size $2n\times2n$.
The question then becomes to bound the probability for a random $W_1\oplus W_2\oplus W_3$ to be stabilized by a non-scalar matrix of the above form. This can be formulated as the following linear algebraic problem. (Recall the correspondence between random $m$-dimensional subspaces and random rank-$m$ matrices as discussed at the beginning of the subsection.) Let $W=\begin{bmatrix}W_{11}&W_{12}&W_{13}\\W_{21}&W_{22}&W_{23}\\W_{31}&W_{32}&W_{33}\end{bmatrix}\in\mathrm{GL}(6n,q)$ be a block matrix where $W_{ij}$ is of size $2n\times2n$. Suppose the columns of $\begin{bmatrix}W_{1i}\\W_{2i}\\W_{3i}\end{bmatrix}$ span $W_i$. Then $D=\operatorname{diag}(D_1,D_2,D_3)$ stabilizes $W_1\oplus W_2\oplus W_3$ if and only if there exists a block diagonal matrix $E=\operatorname{diag}(E_1,E_2,E_3)$, $E_i\in\mathrm{GL}(2n,q)$, such that
\begin{equation}\tag{5}
\begin{bmatrix}D_1&0&0\\0&D_2&0\\0&0&D_3\end{bmatrix}\begin{bmatrix}W_{11}&W_{12}&W_{13}\\W_{21}&W_{22}&W_{23}\\W_{31}&W_{32}&W_{33}\end{bmatrix}=\begin{bmatrix}W_{11}&W_{12}&W_{13}\\W_{21}&W_{22}&W_{23}\\W_{31}&W_{32}&W_{33}\end{bmatrix}\begin{bmatrix}E_1&0&0\\0&E_2&0\\0&0&E_3\end{bmatrix}.
\end{equation}
Note that each direct sum decomposition $W_1\oplus W_2\oplus W_3$, $\dim(W_i)=2n$, has $6\cdot|\mathrm{GL}(2n,q)|^3$ such matrix representations. (The factor 6 takes care of the orders of the three summands.) So the question becomes to bound the probability for a random invertible matrix to have a non-scalar $D$ and $E$ satisfying Equation5.
First, note that Equation5 holds if and only if $D_iW_{i,j}=W_{i,j}E_j$ for $i,j\in[3]$.
\begin{authorclaim}{Claim20}
When $q=\Omega(1)$, we have $\Pr[\forall i,j\in[3],\operatorname{rk}(W_{i,j})=2n]\ge1-20/q$.
\end{authorclaim}
\begin{proof}
Let us work in the setting when $W$ is a random matrix, not necessarily the one above. Then $\Pr[\operatorname{rk}(W)=6n]\ge1-2/q$. For any $i,j\in[3]$, $\Pr[\operatorname{rk}(W_{i,j})<2n]\le2/q$, so $\Pr[\exists i,j\in[3],\operatorname{rk}(W_{i,j})<2n]\le18/q$. It follows that
\begin{align*}
&\Pr[\exists i,j\in[3],\operatorname{rk}(W_{i,j})<2n\mid\operatorname{rk}(W)=6n]\\
&=\frac{\Pr[\exists i,j\in[3],\operatorname{rk}(W_{i,j})<2n\wedge\operatorname{rk}(W)=6n]}{\Pr[\operatorname{rk}(W)=6n]}\\
&\le\frac{18/q}{1-2/q}=\frac{18}{q-2}\le\frac{20}q,
\end{align*}
where the last inequality uses that $q=\Omega(1)$.
\end{proof}
So we assume that $\operatorname{rk}(W_{i,j})=2n$ for all $i,j\in[3]$ in the following, with a loss of probability $\le20/q$.
For $i\in[3]$, by $D_iW_{ii}=W_{ii}E_i$, we have $D_i=W_{ii}E_iW_{ii}^{-1}$. For $i\ne j$, by $(W_{jj}E_jW_{jj}^{-1})W_{ji}=D_jW_{ji}=W_{ji}E_i$, we have $E_j=W_{jj}^{-1}W_{ji}E_iW_{ji}^{-1}W_{jj}$. Again for $i\ne j$, we have $W_{ii}E_iW_{ii}^{-1}W_{ij}=D_iW_{ij}=W_{ij}E_j=W_{ij}W_{jj}^{-1}W_{ji}E_iW_{ji}^{-1}W_{jj}$. It follows that
\[\forall i,j\in[3],i\ne j,\quad E_iW_{ii}^{-1}W_{ij}W_{jj}^{-1}W_{ji}=W_{ii}^{-1}W_{ij}W_{jj}^{-1}W_{ji}E_i.\]
In particular, $E_3$ commutes with $X=W_{33}^{-1}W_{32}W_{22}^{-1}W_{23}$ and $Y=W_{33}^{-1}W_{31}W_{11}^{-1}W_{13}$. Since $W_{ij}$ are independent random invertible matrices, $X$ and $Y$ are independent random invertible matrices. We now resort to the following classical result.
\begin{authorclaim}{Theorem21 ([41], cf. also [23,40])}
Let $X$ and $Y$ be two random matrices in $\mathrm{SL}(n,q)$. Then the probability of $X$ and $Y$ not generating $\mathrm{SL}(n,q)$ is $\le1/q^{\Omega(n)}$.
\end{authorclaim}
Back to our setting, the above theorem implies that the group $G$ generated by random $X$ and $Y$ from $\mathrm{GL}(2n,q)$ contains $\mathrm{SL}(2n,q)$ with probability $\ge1-1/q^{\Omega(n)}$. It follows that $E_3$ belongs to the centralizer of $G$, so $E_3$ must be a scalar matrix. Then note that $D_i$'s and other $E_i$'s are all conjugates of $E_3$. So we have $\forall i\in[3],D_i=E_i=\lambda I_{2n}$ for some $\lambda\ne0\in\mathbb F_q$.
Summarizing the above, we have
\begin{align*}
&\Pr[\mathcal S\text{ non-scalar for }V_1,\ldots,V_6]\\
&\le\Pr[\mathcal S\text{ non-scalar for }V_i\wedge\mathbb F_q^{6n}=U_1\oplus U_2\oplus U_3=W_1\oplus W_2\oplus W_3]+4/q\\
&\le\Pr[\mathcal S\text{ non-scalar for }V_i\mid\mathbb F_q^{6n}=U_1\oplus U_2\oplus U_3=W_1\oplus W_2\oplus W_3]+4/q\\
&\le\Pr[D\text{ non-scalar for }W\wedge\forall i,j\in[3],\operatorname{rk}(W_{ij})=2n]+20/q+4/q\\
&\le\Pr[D\text{ non-scalar for }W\mid\forall i,j\in[3],\operatorname{rk}(W_{ij})=2n]+24/q\\
&\le1/q^{\Omega(n)}+24/q\le1/q^{\Omega(1)},
\end{align*}
when $q=n^{\Omega(1)}$. This concludes the proof of Proposition15. \hfill$\square$
